\documentclass[11pt]{article}

\usepackage[T1]{fontenc}
\usepackage{lmodern}
\usepackage[margin=1.08in]{geometry}
\usepackage{amsmath,amssymb,amsthm,mathtools}
\usepackage{microtype}
\usepackage[backend=biber,style=numeric-comp,sorting=nyt,maxbibnames=99]{biblatex}
\addbibresource{references.bib}
\usepackage[colorlinks=true,linkcolor=blue,citecolor=blue,urlcolor=blue]{hyperref}
\usepackage[nameinlink,capitalise,noabbrev]{cleveref}
\usepackage{enumitem}
\usepackage{booktabs}

\hypersetup{
  pdftitle={Conjugate fibres and irrationality exponents},
  pdfauthor={Christopher D. Long},
  pdfsubject={Irrationality exponents of algebraic exponential fibres}
}
\setlist{topsep=4pt,itemsep=3pt}
\setlength{\emergencystretch}{2em}

\newcommand{\CC}{\mathbb C}
\newcommand{\RR}{\mathbb R}
\newcommand{\QQ}{\mathbb Q}
\newcommand{\ZZ}{\mathbb Z}
\newcommand{\NN}{\mathbb Z_{\ge 0}}
\newcommand{\Qbar}{\overline{\mathbb Q}}
\newcommand{\OO}{\mathcal O}
\newcommand{\eps}{\varepsilon}
\newcommand{\Comp}{\Sigma_{\xi}(a,b)}
\DeclareMathOperator{\Hom}{Hom}
\DeclareMathOperator{\lcm}{lcm}
\DeclareMathOperator{\ord}{ord}
\DeclareMathOperator{\Tr}{Tr}
\DeclareMathOperator{\Nm}{N}

\newtheorem{theorem}{Theorem}[section]
\newtheorem{proposition}[theorem]{Proposition}
\newtheorem{lemma}[theorem]{Lemma}
\newtheorem{corollary}[theorem]{Corollary}
\newtheorem{conjecture}[theorem]{Conjecture}
\theoremstyle{definition}
\newtheorem{definition}[theorem]{Definition}
\theoremstyle{remark}
\newtheorem{remark}[theorem]{Remark}
\numberwithin{equation}{section}

\title{Conjugate fibres and irrationality exponents}
\author{Christopher D. Long\\[2pt]\small \texttt{galizur@gmail.com}}
\date{Draft of October 12, 2026}

\begin{document}
\maketitle

\begin{abstract}
Let $\xi\in\RR^\times$ and let $a,b\in\Qbar^\times$ satisfy
$e^{a\xi}=b$.  Put $F=\QQ(a,b)$, let $d=[F:\QQ]$, and let $r$ be the
number of embeddings $\sigma:F\hookrightarrow\CC$ for which the
conjugate relation
\[
  e^{\sigma(a)\xi}=\sigma(b)
\]
still holds.  We prove
\[
  \mu(\xi)\le \frac{2d}{r}.
\]
The proof combines the separated-weight interpolation determinant of
Liu--Jiang--Zhang with the product formula.  At each compatible
embedding the rational approximation to $\xi$ produces the full
analytic determinant saving; at the remaining embeddings a crude
holomorphic estimate suffices.  We then optimize over all algebraic
exponential certificates, over powers of a fixed certificate, and over
the lattice generated by several certificates.  In particular, full
compatibility after taking a power implies $\mu(\xi)=2$, and for every
positive algebraic $\alpha\ne1$ we obtain
\[
  \mu(\log\alpha)\le
  2\min_{N\ge1}[\QQ(\alpha^N):\QQ].
\]
\end{abstract}

\section{Introduction}

For an irrational real number $x$, its \emph{irrationality exponent} is
\[
 \mu(x)=\sup\left\{\nu>0:
  0<\left|x-\frac pq\right|<q^{-\nu}
  \text{ for infinitely many }p\in\ZZ,\ q\in\ZZ_{>0}\right\}.
\]
Dirichlet's theorem gives $\mu(x)\ge2$ for every irrational $x$; see,
for example, \cite[Chapter~I]{Cassels1957}.  Roth's theorem gives equality
for every irrational algebraic number \cite{Roth1955}.  For explicit
transcendental constants, proving the same exponent is generally much
harder.

A separated-weight interpolation determinant was recently used to prove
$\mu(\pi)=2$ \cite{OpenAI2026}.  Liu, Jiang, and Zhang then developed a
moving-centre form of the interpolation theorem and proved
\[
 \mu(\log q)=2
 \qquad(q\in\QQ_{>0}\setminus\{1\})
\]
\cite{LiuJiangZhang2026}.  Their paper explicitly points out the next
arithmetic obstruction: for an algebraic argument the cleared
determinant is an algebraic integer, so a lower bound at one embedding
must be replaced by estimates for its conjugates.  The purpose of this
note is to carry out that norm argument when a specified collection of
conjugate exponential fibres remains compatible.

Fix $\xi\in\RR^\times$ and $a,b\in\Qbar^\times$ with $e^{a\xi}=b$.
Let
\[
 F=\QQ(a,b),\qquad d=[F:\QQ],\qquad
 \Sigma_F=\Hom_{\QQ}(F,\CC),
\]
and define the set of compatible conjugate fibres by
\begin{equation}\label{eq:compatible-set}
 \Comp=
 \left\{\sigma\in\Sigma_F:
   e^{\sigma(a)\xi}=\sigma(b)\right\}.
\end{equation}
The distinguished inclusion $F\subset\CC$ belongs to $\Comp$, so
$1\le r:=|\Comp|\le d$.

\begin{theorem}[Conjugate-fibre theorem]\label{thm:main}
Let $\xi\in\RR^\times$ and $a,b\in\Qbar^\times$ satisfy $e^{a\xi}=b$.
With $d$ and $r$ as above, for every real number
\[
  \nu>\frac{2d}{r}
\]
there is $Q=Q(a,b,\xi,\nu)$ such that
\[
  \left|\xi-\frac pq\right|\ge q^{-\nu}
  \qquad(p\in\ZZ,\ q\in\ZZ_{>0},\ q\ge Q).
\]
Consequently
\[
  \boxed{\mu(\xi)\le\frac{2d}{r}}.
\]
If every embedding is compatible, so that $r=d$, then
\[
  \boxed{\mu(\xi)=2}.
\]
\end{theorem}

The last assertion uses only Dirichlet's theorem and the fact that
$\xi$ is transcendental.  Indeed, if $\xi$ were a nonzero algebraic
number, then $a\xi$ would be nonzero algebraic, whereas
$e^{a\xi}=b$ is algebraic, contradicting the Hermite--Lindemann theorem
\cite[Chapter~1]{Baker1975}.

The proof has three inputs.  First, the bounded-fibre interpolation
theorem of Liu--Jiang--Zhang produces a nonzero full-row minor for the
centres
\[
  (b^j,ja p_1/q_1,\ldots,ja p_m/q_m),\qquad 0\le j<K.
\]
This covers both distinct multiplicative coordinates and roots of
unity.  Second, at a compatible embedding $\sigma$, the identity
\[
 e^{j\sigma(a)\xi}=\sigma(b)^j
\]
permits the same exact row translation and Taylor-collision estimate as
in \cite{LiuJiangZhang2026}.  Third, after a uniform denominator
clearing, the determinant is a nonzero algebraic integer.  The norm
lower bound therefore compares the strong estimates at the $r$
compatible embeddings with crude growth estimates at the other $d-r$
embeddings.  The two leading terms are respectively $r\nu$ and $d$;
the sharp parameter threshold for this scalar norm construction is
$2d/r$.

Sections~\ref{sec:interpolation}--\ref{sec:proof} prove
\cref{thm:main}.  Section~\ref{sec:saturation} gives the optimization,
power-saturation, torsion-defect, and relation-lattice consequences.

\section{The interpolation input}\label{sec:interpolation}

We recall the precise form of the geometric input that is needed.  The following is \cite[Theorem~9]{LiuJiangZhang2026}.  Fix
positive rational weights
\[
 W=(w_0,w_1,\ldots,w_m),\qquad
 V=(v_0,w_1/\theta,\ldots,w_m/\theta),
 \qquad 0<\theta<1.
\]
Let $\mathcal P_W(H)$ be the span of the monomials $Y^hX^\alpha$ with
\[
 w_0h+\sum_{i=1}^m w_i\alpha_i\le H,
\]
and let $\mathcal J_V(H)$ be the quotient of
$\CC[[t,u_1,\ldots,u_m]]$ retaining precisely the monomials
$t^su^\beta$ with
\[
 v_0s+\frac1\theta\sum_{i=1}^m w_i\beta_i<H.
\]

\begin{theorem}[Bounded-fibre moving-centre interpolation]\label{thm:MI}
Fix $m,K\ge1$, $1\le\kappa\le K$, positive rational $w_0,v_0$, and
$0<\theta<1$ such that
\begin{equation}\label{eq:MI-volume}
 K\left(\frac{w_0}{v_0}\right)\theta^m<1,
 \qquad \kappa\theta^m<1.
\end{equation}
There are successive lower thresholds for the positive rational weights
$w_1,\ldots,w_m$, independent of the interpolation centres, with the
following property.  Let
\[
 (y_j,c_{j1},\ldots,c_{jm})\in\CC^\times\times\CC^m,
 \qquad 0\le j<K,
\]
assume that at most $\kappa$ indices have a given value of $y_j$, and
assume that
\begin{equation}\label{eq:fibre-separation}
 j\ne\ell,\ y_j=y_\ell
 \quad\Longrightarrow\quad
 c_{ji}\ne c_{\ell i}\quad(1\le i\le m).
\end{equation}
Then, for every sufficiently large admissible degree $H$, the map
\[
 \mathcal P_W(H)\longrightarrow
 \bigoplus_{j=0}^{K-1}\mathcal J_V(H)
\]
obtained by substituting
\[
 Y=y_j(1+t),\qquad
 X_i=c_{ji}+u_i+\log(1+t)
\]
is onto.  The divisibility modulus for $H$ depends only on the weights;
the eventual threshold may depend on the centres.
\end{theorem}

The proof in \cite{LiuJiangZhang2026} uses a local intersection-length
estimate, separated weights, a curve inequality, an ordinary blow-up,
and Serre vanishing.  We shall also use the determinant translation and
collision estimates from Sections~4--6 of that paper.  Those estimates
are uniform over the selected full-row minor and its columns.

\subsection{Parameters}

The following elementary lemma isolates the only change from the
parameter choice for a single rational logarithm.

\begin{lemma}[Gap parameters]\label{lem:parameters}
Let $d,r$ be positive integers and let $\nu>2d/r$.  There exist rational
numbers
\[
 0<\theta<A<B<1,\qquad 0<\eta<1,
\]
and rational numbers $C>1$ such that
\begin{align}
 &A^2<\theta<A,
 \label{eq:param1}\\
 &g:=r\nu\bigl(A(1-\eta)-\theta\bigr)-d(1-\theta)>0,
 \label{eq:param-gap}\\
 &\frac A\theta<C<\frac1A,
 \qquad
 A<B<\min\{C^{-1},C\theta\}.
 \label{eq:param-CB}
\end{align}
In particular $C\theta<1$, $CB<1$, $C\theta/B>1$, and $B/A>1$.
\end{lemma}

\begin{proof}
Choose a rational number
\[
 \frac12<\rho<1-\frac{d}{r\nu};
\]
this interval is nonempty because $r\nu>2d$.  For a sufficiently small
positive rational $\delta$, put
\[
 \theta=1-\delta,\qquad A=1-\rho\delta.
\]
Then $A>\theta$ and
\[
 \theta-A^2=(2\rho-1)\delta-\rho^2\delta^2>0.
\]
Moreover
\[
 r\nu(A-\theta)-d(1-\theta)
 =\delta\bigl(r\nu(1-\rho)-d\bigr)>0.
\]
Choose a sufficiently small positive rational $\eta$ to preserve this
strict inequality after replacing $A$ by $A(1-\eta)$.  Since
$A^2<\theta$, the interval $(A/\theta,1/A)$ is nonempty; choose rational
$C$ in it.  Both $C^{-1}$ and $C\theta$ exceed $A$, so a rational $B$
satisfying \eqref{eq:param-CB} exists.  Finally
$C\theta<\theta/A<1$.
\end{proof}

For a positive integer $m$, set
\begin{equation}\label{eq:dimension-parameters}
 K=\lfloor C^m\rfloor,
 \qquad w_0=B^{-m},
 \qquad v_0=2K\theta^m w_0.
\end{equation}
Then
\begin{equation}\label{eq:parameter-limits}
 K\left(\frac{w_0}{v_0}\right)\theta^m=\frac12,
 \qquad K\theta^m\longrightarrow0,
 \qquad \frac K{w_0}\longrightarrow0,
 \qquad v_0\longrightarrow\infty.
\end{equation}
For later use put
\begin{equation}\label{eq:L-def}
 L=\frac{\eta^2K\theta^m}{(m+1)v_0A^m}
   =\frac{\eta^2}{2(m+1)}\left(\frac BA\right)^m.
\end{equation}
Thus $L\to\infty$ as $m\to\infty$.

\section{The number-field determinant}\label{sec:determinant}

Fix the data of \cref{thm:main}, and suppose for contradiction that
there is a real $\nu>2d/r$ for which
\begin{equation}\label{eq:good-approximations}
 \left|\xi-\frac pq\right|\le q^{-\nu}
\end{equation}
holds for arbitrarily large positive integers $q$.  Replacing $\nu$ by
a slightly smaller rational number still larger than $2d/r$, we may
assume that $\nu$ is rational.

Choose the parameters of \cref{lem:parameters}.  Fix a rational
truncation parameter
\begin{equation}\label{eq:Phi}
 \Phi>\frac2\theta.
\end{equation}
We shall first take $m$ large, then choose $m$ approximations
\[
 x_i=\frac{p_i}{q_i},\qquad
 \left|\xi-x_i\right|\le q_i^{-\nu},
 \qquad w_i=\lceil\log q_i\rceil,
 \qquad 1\le i\le m,
\]
successively.  The denominators are taken large enough that $x_i\ne0$,
$|x_i|\le|\xi|+1$, and the weights meet the successive thresholds in
\cref{thm:MI}.  This is possible because those thresholds are
independent of the centres.

Use the centres
\begin{equation}\label{eq:centres}
 y_j=b^j,\qquad c_{ji}=ja x_i,
 \qquad 0\le j<K,\quad1\le i\le m.
\end{equation}
We apply \cref{thm:MI} with $\kappa=K$.  The second inequality in
\eqref{eq:MI-volume} holds for large $m$ by $K\theta^m\to0$.  If
$y_j=y_\ell$ and $j\ne\ell$, then
\[
 c_{ji}-c_{\ell i}=(j-\ell)a x_i\ne0
\]
for every $i$, so \eqref{eq:fibre-separation} holds.  Thus the same
argument covers roots of unity and non-torsion values of $b$.

Set
\begin{equation}\label{eq:truncations}
 T_i=\left\lceil\frac{\Phi w_i}{v_0}\right\rceil,
 \qquad
 G_i(t)=\sum_{1\le k<T_i}\frac{(-1)^{k+1}}{k}t^k,
 \qquad
 L_i=\lcm(1,\ldots,T_i-1),
\end{equation}
where an empty sum and empty least common multiple are interpreted as
$0$ and $1$.  Because $v_0T_i>w_i/\theta$, replacing $\log(1+t)$ by
$G_i(t)$ is an invertible filtered change on every retained jet packet
\cite[Section~4]{LiuJiangZhang2026}.

For an admissible degree $H$, use the monomials
\[
 P=Y^hX^\alpha,
 \qquad w_0h+\sum_iw_i\alpha_i\le H,
\]
as columns, and use the triples
\[
 \varrho=(j,s,\beta),\qquad
 0\le j<K,
 \qquad v_0s+\frac1\theta\sum_iw_i\beta_i<H,
\]
as rows.  The matrix entry is
\begin{equation}\label{eq:matrix-entry}
 [t^su^\beta]\,
 b^{jh}(1+t)^h
 \prod_{i=1}^m(ja x_i+G_i(t)+u_i)^{\alpha_i}.
\end{equation}
For every sufficiently large admissible $H$, interpolation supplies a
nonzero square minor $\Delta_H\in F^\times$ containing every row.  Let
$M=M_H$ be its size, and put
\begin{equation}\label{eq:bbar}
 \bar\beta_H=\frac1{MH}
 \sum_{\varrho=(j,s,\beta)}\sum_iw_i\beta_i.
\end{equation}
Then
\begin{equation}\label{eq:row-asymptotic}
 M\sim
 \frac{K\theta^mH^{m+1}}
 {(m+1)!v_0\prod_iw_i},
 \qquad 0\le\bar\beta_H\le\theta.
\end{equation}

\subsection{Algebraic-integral clearing and the norm inequality}

Choose an integer $D\ge1$ such that
\[
 Da,\ Db\in\OO_F.
\]
Put
\[
 \mathfrak D_H=\prod_{i=1}^mL_i^{\lfloor H/w_i\rfloor}.
\]
Scale the column $(h,\alpha)$ and the row $(j,s,\beta)$ by
\begin{equation}\label{eq:scales}
 C_{h,\alpha}=D^{(K-1)h}\prod_i(Dq_i)^{\alpha_i},
 \qquad
 S_{j,s,\beta}=\mathfrak D_H\prod_i(Dq_i)^{-\beta_i}.
\end{equation}

\begin{lemma}[Norm lower bound]\label{lem:norm}
Let $\Lambda=4\log2$ and $w_*=\min_iw_i$.  Define
\begin{equation}\label{eq:Ear}
 E_{\mathrm{ar}}=
 \frac{\Lambda\Phi m}{v_0}
 +\Lambda\sum_{i=1}^m\frac1{w_i}
 +\frac{\theta+\log D}{w_*}
 +\frac{(K-1)\log D}{w_0}.
\end{equation}
Then
\begin{equation}\label{eq:norm-inequality}
 0\le
 \sum_{\sigma\in\Sigma_F}
 \frac{\log|\sigma(\Delta_H)|}{MH}
 +d\bigl(1-\bar\beta_H+E_{\mathrm{ar}}\bigr).
\end{equation}
\end{lemma}

\begin{proof}
The entry obtained from \eqref{eq:matrix-entry} after the row and column
scalings is zero unless $\beta\le\alpha$.  Otherwise it equals
\begin{align*}
 &\mathfrak D_H(Db)^{jh}D^{(K-1-j)h}
 \binom\alpha\beta[t^s](1+t)^h\\
 &\hspace{25mm}\cdot
 \prod_{i=1}^m
 \bigl(j(Da)p_i+Dq_iG_i(t)\bigr)^{\alpha_i-\beta_i}.
\end{align*}
The denominator contributed by the $i$th truncated logarithm divides
$L_i^{\alpha_i-\beta_i}$, and
$\alpha_i-\beta_i\le\lfloor H/w_i\rfloor$.  Hence every scaled entry
belongs to $\OO_F$.  The scaled determinant
\[
 \widetilde\Delta_H
 =\Delta_H\prod_{\text{columns}}C_{h,\alpha}
             \prod_{\text{rows}}S_{j,s,\beta}
\]
is therefore a nonzero algebraic integer.  Consequently
\[
 1\le|\Nm_{F/\QQ}(\widetilde\Delta_H)|
 =\prod_{\sigma\in\Sigma_F}|\sigma(\widetilde\Delta_H)|;
\]
we count complex embeddings separately, as in the usual product formula
\cite[Chapter~I]{Neukirch1999}.

It remains to bound the positive rational scaling factor.  As in
\cite[Section~4]{LiuJiangZhang2026},
\[
 \frac{\log\mathfrak D_H}{H}
 \le\frac{\Lambda\Phi m}{v_0}
    +\Lambda\sum_i\frac1{w_i}.
\]
Every selected column satisfies
\[
 \sum_i\alpha_i\log q_i\le H,
 \qquad |\alpha|\le\frac H{w_*},
 \qquad h\le\frac H{w_0}.
\]
Moreover, since $0\le w_i-\log q_i<1$,
\[
 \sum_{\text{rows}}\sum_i\beta_i\log q_i
 \ge MH\bar\beta_H-\frac{MH\theta}{w_*}.
\]
The negative row contribution involving $\log D$ may be discarded.
Thus the logarithm of the total scaling factor, divided by $MH$, is at
most $1-\bar\beta_H+E_{\mathrm{ar}}$.  Taking logarithms of the norm
inequality proves \eqref{eq:norm-inequality}.
\end{proof}

\section{Archimedean determinant estimates}\label{sec:analytic}

Set
\[
 C_a=\max_{\sigma\in\Sigma_F}\max\{1,|\sigma(a)|\},
 \qquad
 C_b=\max_{\sigma\in\Sigma_F}\max\{1,|\sigma(b)|\},
 \qquad X=|\xi|+1.
\]
All constants in this section are uniform over the finite set
$\Sigma_F$.

\subsection{Compatible embeddings}

For $\sigma\in\Comp$, define
\[
 z_j(t)=j\sigma(a)\xi+\log(1+t),
 \qquad
 \epsilon_{ji}=j\sigma(a)(x_i-\xi),
 \qquad
 \tau_i(t)=G_i(t)-\log(1+t).
\]
Compatibility gives the exact identities
\begin{equation}\label{eq:compatible-identities}
 e^{z_j(t)}=\sigma(b)^j(1+t),
 \qquad
 j\sigma(a)x_i+G_i(t)+u_i
 =z_j(t)+u_i+\epsilon_{ji}+\tau_i(t).
\end{equation}
For a selected column $P=Y^hX^\alpha$ and a multi-index $\lambda$, let
\[
 f_{\lambda,P}(z)
 =[u^\lambda]P(e^z,z+u_1,\ldots,z+u_m)
 =\binom\alpha\lambda e^{hz}z^{|\alpha|-|\lambda|},
\]
with the convention that this is zero unless $\lambda\le\alpha$.
Expanding \eqref{eq:compatible-identities} gives the same exact row
translation as in \cite[Section~5.1]{LiuJiangZhang2026}:
\begin{align}\label{eq:row-translation}
 \operatorname{row}_{j,s,\beta}^{\sigma}
 ={}&\sum_{\substack{\lambda\ge\beta\\w\lambda\le H}}
 \binom\lambda\beta
 \sum_{0\le e\le\lambda-\beta}
 \binom{\lambda-\beta}{e}
 \epsilon_j^{\lambda-\beta-e}\\
 &\quad\cdot\sum_{k=0}^s
 [t^k]\prod_i\tau_i(t)^{e_i}
 \bigl([t^{s-k}]f_{\lambda,P}(z_j(t))\bigr)_P.
 \nonumber
\end{align}
Here and below $w\lambda=\sum_iw_i\lambda_i$.

Since
\[
 |\epsilon_{ji}|
 \le KC_aq_i^{-\nu}
 \le KC_ae^{\nu}e^{-\nu w_i},
\]
and a nonzero tail coefficient in \eqref{eq:row-translation} uses
$we<H/\Phi$, every scalar multiplying a translated test row is bounded
by
\begin{equation}\label{eq:translation-scalar}
 \exp\{-\nu w(\lambda-\beta)+HE_{\mathrm{tr}}\},
\end{equation}
where one may take
\begin{equation}\label{eq:Etr}
 E_{\mathrm{tr}}
 =\frac\nu\Phi+\frac{\log2}{v_0}
 +\frac{\log4+\log(2KC_a)+\nu}{w_*}.
\end{equation}

Choose
\[
 R=100K(1+C_a|\xi|).
\]
On $|t|\le1/2$, every $z_j(t)$ lies in $|z|<R/2$.  The Taylor-collision
argument of \cite[Section~5.2]{LiuJiangZhang2026}, originally appearing
in this interpolation-determinant form in \cite{Laurent2000}, gives the
following uniform estimate.  Put
\begin{equation}\label{eq:Ehol}
 E_{\mathrm{hol}}
 =\frac R{w_0}+\frac{\log2}{v_0}
  +\frac{\log(2R)}{w_*},
 \qquad c=\frac{\log2}{4}.
\end{equation}
Let
\[
 S=\sum_iw_i,
 \qquad
 L_H=L\left(1+\frac{S}{AH}\right)^{-m}.
\]
Thus $L_H\to L$ as $H\to\infty$.

\begin{lemma}[Compatible-place estimate]\label{lem:compatible}
There is a quantity $\delta_H\to0$, independent of
$\sigma\in\Comp$ and of the selected full-row minor, such that
\begin{equation}\label{eq:compatible-bound}
 \frac{\log|\sigma(\Delta_H)|}{MH}
 \le E_{\mathrm{tr}}+E_{\mathrm{hol}}+\delta_H
 +\max\left\{-cL_H,
 -\nu\bigl(A(1-\eta)-\bar\beta_H\bigr)\right\}.
\end{equation}
\end{lemma}

\begin{proof}
We indicate the complete reduction to the cited collision estimate.
The row expansion \eqref{eq:row-translation} has only polynomially many
terms in $H$ for each row.  Choose one translated row from each original
row and let $n_\lambda$ be the number chosen with transverse index
$\lambda$.  Expanding the entire functions $f_{\lambda,P}$ in Taylor
series in $z/R$, any nonzero determinant contribution uses distinct
Taylor degrees within each fixed $\lambda$-group.  Hence
\[
 |\det T|
 \le\exp\left\{-c\sum_\lambda n_\lambda^2
 +MH(E_{\mathrm{hol}}+o(1))\right\}.
\]
This is exactly the collision estimate of
\cite[Section~5.2]{LiuJiangZhang2026}; its proof is uniform in the
columns and depends on the fibre only through the radius $R$.

Let
\[
 D_A(H)=\#\{\lambda\in\NN^m:w\lambda\le AH\}.
\]
The same unit-cube estimates as in \cite[Section~6]{LiuJiangZhang2026}
give
\[
 M\ge\frac{K\theta^mH^{m+1}}
 {(m+1)!v_0\prod_iw_i},
 \qquad
 D_A(H)\le\frac{(AH+S)^m}{m!\prod_iw_i}.
\]
If at least $\eta M$ selected indices have weight at most $AH$, then
Cauchy--Schwarz yields
\[
 \frac1{MH}\sum_\lambda n_\lambda^2\ge L_H.
\]
Otherwise more than $(1-\eta)M$ indices have weight greater than $AH$.
Combining this with the scalar estimate
\eqref{eq:translation-scalar} gives the exponent
$-\nu(A(1-\eta)-\bar\beta_H)$.  The number of multilinear expansion
terms is polynomial to the power $M$; after division by $MH$ its
contribution is $o(1)$.  This proves \eqref{eq:compatible-bound}.
\end{proof}

\subsection{The remaining embeddings}

At an incompatible embedding no approximation gain is available, but
none is needed.  A direct Cauchy estimate gives a normalized cost that
can be made arbitrarily small by the separated parameter choices.

\begin{lemma}[Crude estimate at every embedding]\label{lem:crude}
For every $\sigma\in\Sigma_F$,
\begin{equation}\label{eq:crude-bound}
 \frac{\log|\sigma(\Delta_H)|}{MH}
 \le E_{\mathrm{bad}}+\frac{\log M}{2H},
\end{equation}
where
\begin{equation}\label{eq:Ebad}
 E_{\mathrm{bad}}
 =\frac{K\log C_b+\log(3/2)}{w_0}
 +\frac{\log\bigl(2(1+KC_aX+\log2)\bigr)}{w_*}
 +\frac{\log2}{v_0}.
\end{equation}
\end{lemma}

\begin{proof}
The conjugate of \eqref{eq:matrix-entry} is zero unless
$\beta\le\alpha$.  Otherwise it is
\[
 \sigma(b)^{jh}\binom\alpha\beta[t^s](1+t)^h
 \prod_i(j\sigma(a)x_i+G_i(t))^{\alpha_i-\beta_i}.
\]
On $|t|=1/2$ we have $|G_i(t)|\le\log2$, while
$|x_i|\le X$, $j<K$, $s\le H/v_0$, $h\le H/w_0$, and
$|\alpha|\le H/w_*$.  Cauchy's estimate and
$\binom\alpha\beta\le2^{|\alpha|}$ therefore bound each entry by
$\exp(HE_{\mathrm{bad}})$.  Hadamard's inequality contributes at most
$M^{M/2}$, which gives \eqref{eq:crude-bound}.
\end{proof}

\subsection{Making the errors small}

The choices are made in the following order.  First choose $\Phi$ so
large that the contribution $r\nu/\Phi$ is smaller than any prescribed
positive tolerance.  Next let $m$ tend to infinity.  By
\eqref{eq:parameter-limits} and \eqref{eq:L-def},
\[
 \frac{m}{v_0}\to0,
 \qquad \frac K{w_0}\to0,
 \qquad \frac1{v_0}\to0,
 \qquad L\to\infty.
\]
The constants $D,C_a,C_b,\xi,d,r$ are fixed.  Thus $m$ may be chosen so
that all terms in \eqref{eq:Ear}, \eqref{eq:Ehol}, and
\eqref{eq:Ebad} not involving $w_i$ are as small as desired, and also
so that
\begin{equation}\label{eq:low-gap}
 rcL>d+1.
\end{equation}
Finally, because the interpolation thresholds are centre-independent,
the approximations may be selected successively with all $w_i$ so
large that
\[
 \sum_i\frac1{w_i},\qquad \frac{1+\log K}{w_*}
\]
are as small as desired.  Once all parameters and centres are fixed,
$H$ tends to infinity through admissible degrees; then
$L_H\to L$, $\delta_H\to0$, and $(\log M)/H\to0$.

\section{Proof of the conjugate-fibre theorem}\label{sec:proof}

Apply \cref{lem:compatible} at the $r$ embeddings in $\Comp$ and
\cref{lem:crude} at the other $d-r$ embeddings.  Substitution in
\eqref{eq:norm-inequality} gives
\begin{align}\label{eq:final-master}
 0\le{}&d(1-\bar\beta_H)
 +r\max\left\{-cL_H,
 -\nu\bigl(A(1-\eta)-\bar\beta_H\bigr)\right\}
 +\mathcal E_H,
\end{align}
where, by the choice order above, the nonnegative error
$\mathcal E_H$ can be made arbitrarily small and then kept fixed while
$H\to\infty$.

For the collision alternative,
\[
 d(1-\bar\beta_H)-rcL_H\le d-rcL_H,
\]
which is strictly negative for all sufficiently large $H$ by
\eqref{eq:low-gap}.  For the large-index alternative, since
$r\nu>d$ and $0\le\bar\beta_H\le\theta$,
\begin{align*}
 &r\nu\bigl(A(1-\eta)-\bar\beta_H\bigr)
   -d(1-\bar\beta_H)\\
 &\qquad\ge
 r\nu\bigl(A(1-\eta)-\theta\bigr)-d(1-\theta)
 =g>0.
\end{align*}
This is exactly the positive gap displayed in \eqref{eq:param-gap}.  Thus both entries of the maximum in
\eqref{eq:final-master}, after adding $d(1-\bar\beta_H)$, are bounded
above by a fixed negative number.  Taking all fixed errors sufficiently
small and then $H$ sufficiently large contradicts
\eqref{eq:final-master}.

Therefore \eqref{eq:good-approximations} cannot hold for arbitrarily
large denominators when $\nu>2d/r$.  This proves the asserted eventual
lower bound and hence $\mu(\xi)\le2d/r$.  As noted in the introduction,
$\xi$ is transcendental, and if $r=d$, Dirichlet's theorem supplies the
opposite inequality $\mu(\xi)\ge2$.  This completes the proof of
\cref{thm:main}.

\begin{remark}[Why the scalar norm threshold is $2d/r$]
The constraint $A^2<\theta<A<1$ implies
\[
 1-\theta-2(A-\theta)=1-2A+\theta>(1-A)^2>0.
\]
Hence
\[
 \frac{1-\theta}{A-\theta}>2.
\]
The leading high-index inequality
$r\nu(A-\theta)>d(1-\theta)$ therefore forces
$\nu>2d/r$.  Conversely, \cref{lem:parameters} constructs admissible
parameters for every strict inequality $\nu>2d/r$.  Thus $2d/r$ is the
sharp threshold of this scalar product-formula architecture, although
it need not be the true irrationality exponent when $r<d$.
\end{remark}

\begin{remark}[Effectivity]
The theorem is qualitative.  As in \cite{LiuJiangZhang2026}, after the
centres have been selected the analytic contradiction has an explicit
large-degree threshold, but the centre-dependent Serre-vanishing
threshold in the interpolation theorem is not made effective.
\end{remark}

\section{Saturation and optimization}\label{sec:saturation}

The certificate $e^{a\xi}=b$ is not canonical.  Powers, products, and
integer combinations may enlarge the proportion of compatible
embeddings or reduce the field degree.  The following corollaries make
that optimization intrinsic.

\subsection{Optimization over all certificates}

Let
\[
 \mathcal C_{\exp}(\xi)=
 \{(a,b)\in\Qbar^\times\times\Qbar^\times:e^{a\xi}=b\}.
\]
For $(a,b)\in\mathcal C_{\exp}(\xi)$, write
\[
 \rho_\xi(a,b)=
 \frac{|\Sigma_\xi(a,b)|}{[\QQ(a,b):\QQ]}.
\]
If $\mathcal C_{\exp}(\xi)\ne\varnothing$, define the
\emph{exponential conjugate-fibre density}
\begin{equation}\label{eq:rhoexp}
 \rho_{\exp}(\xi)=
 \sup_{(a,b)\in\mathcal C_{\exp}(\xi)}\rho_\xi(a,b).
\end{equation}

\begin{corollary}[Optimized conjugate-fibre bound]\label{cor:optimized}
If $\mathcal C_{\exp}(\xi)\ne\varnothing$, then
\[
 \boxed{\mu(\xi)\le\frac{2}{\rho_{\exp}(\xi)}}.
\]
In particular, if $\rho_{\exp}(\xi)=1$, then $\mu(\xi)=2$.
\end{corollary}

\begin{proof}
Apply \cref{thm:main} to every certificate and take the infimum of the
bounds $2/\rho_\xi(a,b)$.
\end{proof}

\subsection{Power saturation}

Fix one certificate $e^{a\xi}=b$.  For $N\ge1$, put
\begin{equation}\label{eq:power-data}
 F_N=\QQ(a,b^N),\qquad d_N=[F_N:\QQ],
\end{equation}
and
\begin{equation}\label{eq:power-r}
 r_N=\#\left\{\sigma:F_N\hookrightarrow\CC:
 e^{N\sigma(a)\xi}=\sigma(b^N)\right\}.
\end{equation}

\begin{corollary}[Power saturation]\label{cor:powers}
For every certificate $e^{a\xi}=b$,
\[
 \boxed{
 \mu(\xi)\le
 2\inf_{N\ge1}\frac{d_N}{r_N}
 =\frac{2}{\sup_{N\ge1}(r_N/d_N)}.}
\]
If $r_N=d_N$ for some $N$, then $\mu(\xi)=2$.
\end{corollary}

\begin{proof}
For every $N\ge1$,
\[
 e^{(Na)\xi}=b^N,
\]
and $\QQ(Na,b^N)=F_N$.  Apply \cref{thm:main} and optimize over $N$.
\end{proof}

Powering may help in two ways: it can identify conjugates of $b$ that
differ only by torsion, and it can replace $\QQ(a,b)$ by the smaller
field $\QQ(a,b^N)$.

\begin{corollary}[Torsion-defect criterion]\label{cor:torsion}
Let $L$ be a normal closure of $\QQ(a,b)$ over $\QQ$.  Suppose that for
every embedding $\tau:L\hookrightarrow\CC$ the defect
\begin{equation}\label{eq:defect}
 \delta_\tau=\frac{e^{\tau(a)\xi}}{\tau(b)}
\end{equation}
is a root of unity.  Then
\[
 \boxed{\mu(\xi)=2}.
\]
\end{corollary}

\begin{proof}
Choose $N$ divisible by the orders of the finitely many roots of unity
$\delta_\tau$.  Then
\[
 e^{N\tau(a)\xi}=\tau(b)^N=\tau(b^N)
\]
for every $\tau$.  Every embedding of $F_N=\QQ(a,b^N)$ extends to an
embedding of $L$, so $r_N=d_N$.  Apply \cref{cor:powers}.
\end{proof}

\subsection{The relation lattice}

Suppose that
\[
 e^{a_j\xi}=b_j,
 \qquad a_j,b_j\in\Qbar^\times,
 \qquad 1\le j\le s.
\]
For $n=(n_1,\ldots,n_s)\in\ZZ^s$, set
\begin{equation}\label{eq:relation-combination}
 A_n=\sum_{j=1}^s n_ja_j,
 \qquad
 B_n=\prod_{j=1}^s b_j^{n_j}.
\end{equation}
Whenever $A_n\ne0$, we have $e^{A_n\xi}=B_n$.  Put
\[
 F_n=\QQ(A_n,B_n),\qquad d_n=[F_n:\QQ],
\]
and
\[
 r_n=\#\left\{\sigma:F_n\hookrightarrow\CC:
 e^{\sigma(A_n)\xi}=\sigma(B_n)\right\}.
\]

\begin{corollary}[Relation-lattice saturation]\label{cor:relations}
With the notation above,
\[
 \boxed{
 \mu(\xi)\le
 2\inf_{\substack{n\in\ZZ^s\\A_n\ne0}}
 \frac{d_n}{r_n}.}
\]
In particular, an integer combination whose conjugate fibres are all
compatible proves $\mu(\xi)=2$, even when none of the original
certificates has full compatibility.
\end{corollary}

\begin{proof}
Apply \cref{thm:main} to the certificate
$e^{A_n\xi}=B_n$ and optimize over $n$.
\end{proof}

In a common normal field, define the individual defects
\[
 \delta_{\tau,j}=\frac{e^{\tau(a_j)\xi}}{\tau(b_j)}.
\]
The combined certificate is compatible at $\tau$ precisely when
\[
 \prod_{j=1}^s\delta_{\tau,j}^{n_j}=1.
\]
When the defects are roots of unity, the search for a fully compatible
combination is therefore a finite system of linear congruences.

\subsection{Positive algebraic logarithms}

For $\alpha\in\Qbar_{>0}\setminus\{1\}$, define its \emph{power degree}
by
\begin{equation}\label{eq:power-degree}
 d_{\mathrm{pow}}(\alpha)=
 \min_{N\ge1}[\QQ(\alpha^N):\QQ].
\end{equation}

\begin{corollary}[Algebraic logarithms]\label{cor:alglog}
For every positive algebraic $\alpha\ne1$,
\[
 \boxed{\mu(\log\alpha)\le2d_{\mathrm{pow}}(\alpha).}
\]
If $\alpha^N\in\QQ$ for some $N\ge1$, then
\[
 \boxed{\mu(\log\alpha)=2}.
\]
\end{corollary}

\begin{proof}
Use the powered certificate
\[
 e^{N\log\alpha}=\alpha^N.
\]
In the field $\QQ(\alpha^N)$, exactly one embedding fixes the generator
$\alpha^N$, so $r_N=1$.  Thus \cref{cor:powers} gives
\[
 \mu(\log\alpha)
 \le2[\QQ(\alpha^N):\QQ]
\]
for every $N$.  Minimize over $N$.  If one power is rational, the upper
bound is $2$, and Dirichlet gives equality.  The rational-power case is
also obtained directly in \cite[Corollary~7]{LiuJiangZhang2026}.
\end{proof}

\section{Examples}\label{sec:examples}

\begin{corollary}\label{cor:pi-examples}
For every positive rational $D$,
\[
 \mu(\pi\sqrt D)=\mu\left(\frac\pi{\sqrt D}\right)=2.
\]
\end{corollary}

\begin{proof}
For $\xi=\pi\sqrt D$, use $a=i/\sqrt D$ and $b=-1$; for
$\xi=\pi/\sqrt D$, use $a=i\sqrt D$ and $b=-1$.  Every embedding sends
$a$ to $a$ or $-a$, and $e^{\pm i\pi}=-1$, so all fibres are compatible.
\end{proof}

Taking $D=1$ recovers $\mu(\pi)=2$, the theorem of
\cite{OpenAI2026}.

\begin{corollary}\label{cor:arctan}
For every nonzero rational number $t$,
\[
 \boxed{\mu(\arctan t)=2}.
\]
\end{corollary}

\begin{proof}
Let $\xi=\arctan t$.  Then
\[
 e^{2i\xi}=\frac{1+it}{1-it}\in\QQ(i).
\]
Complex conjugation sends $2i$ to $-2i$ and the right-hand side to its
inverse, so both embeddings of $\QQ(i)$ are compatible.
\end{proof}

\begin{corollary}[Quadratic units]\label{cor:quadratic-units}
Let $K=\QQ(\sqrt D)$ be a real quadratic field, and let
$\varepsilon>0$ be a unit.  If $\Nm_{K/\QQ}(\varepsilon)=1$, then
\[
 \mu\left(\frac{\log\varepsilon}{\sqrt D}\right)=2.
\]
The same conclusion holds for norm $-1$ after replacing the certificate
by its square.
\end{corollary}

\begin{proof}
Put $\xi=\log\varepsilon/\sqrt D$, $a=\sqrt D$, and $b=\varepsilon$.
For norm $1$, the nontrivial automorphism sends
$a\mapsto-a$ and $b\mapsto b^{-1}$, so both fibres are compatible.  For
norm $-1$, square the relation and apply \cref{cor:powers}.
\end{proof}

\section*{Acknowledgment of inputs}

The geometric interpolation theorem and the analytic row-translation
and collision estimates used here are those of Liu--Jiang--Zhang
\cite{LiuJiangZhang2026}, building on the interpolation-determinantmethod of Laurent \cite{Laurent1991,Laurent2000} and the recent proof for
$\pi$ \cite{OpenAI2026}.  The new step in this note is the simultaneous
number-field denominator clearing and product-formula comparison across
compatible and incompatible conjugate fibres.

\printbibliography

\end{document}